\section{A natural-density improvement over the greedy ideals in every number field}
\label{sec:ideal-natural-gain}

Best, Huan, McNew, Miller, Powell, Tor and Weinstein ask whether the
greedy progression-free integral ideals in a quadratic number field can
be surpassed in natural density
(\href{https://arxiv.org/pdf/1412.0999}{Question~5.1 of their paper}).
The answer is affirmative in every number field. The argument extends
the finite valuation-cube exchange in the earlier
\href{https://github.com/VladimirReshetnikov/ProveIt/blob/5cc267af75a635199ac5d93a7d49cbce521daefc/docs/incoming/grh-progression-densities-and-cubic-prime-race.zip}{ProveIt polynomial-progression report}.
Its integer precedent is McNew's exchange construction. The new step
is to prove that a profitable finite exchange always exists over an
arbitrary number field, independently of its splitting patterns. The
proof uses no RH or GRH assumption. Historical priority is not asserted.

Let $K$ be a number field of degree $d$, let $\mathcal I_K$ be its nonzero
integral ideals, and put
\[
 \begin{aligned}
 A&=\left\{\sum_{j\ge0}\epsilon_j3^j:\epsilon_j\in\{0,1\},
                         \ \epsilon_j=0\text{ eventually}\right\},\\
 \mathcal G_K&=\{I\in\mathcal I_K:v_{\mathfrak p}(I)\in A
                         \text{ for every prime ideal }\mathfrak p\}.
 \end{aligned}
\]
Natural density is relative to the ideals of norm at most $X$.
A forbidden progression in this section is
$I,IJ,IJ^2$ with $J$ a proper nonzero \emph{integral} ideal.

\begin{theorem}\label{thm:ideal-natural-gain}
For every number field $K$, an explicitly decidable set
$\mathcal E\subseteq\mathcal I_K$ avoids all such progressions, has
natural density, and satisfies
\[
                         d(\mathcal E)>d(\mathcal G_K).
\]
There is an effective search for a finite set of prime ideals specifying
$\mathcal E$, with a positive rational certificate for the density gain.
\end{theorem}

\paragraph{The finite exchange.}
For a nonempty finite set $S$ of prime ideals, write $k=|S|$ and replace
the permitted local vectors $A^k$ by
\[
 D_S=\{u\in A^k:\min_i u_i=0\}
       \ \cup\ \bigl(\{1,2\}^k\setminus\{(1,\ldots,1)\}\bigr).
\]
Let $\mathcal E_S$ require this condition at $S$ and retain the condition
$v_{\mathfrak p}(I)\in A$ at every prime ideal outside $S$.

The set $A$ contains no nonconstant three-term arithmetic progression.
Indeed, in $a+c=2b$ the ternary digit sums have no carries, and a digit of
$2b$ is either zero or two. Equality forces the corresponding digits
of $a,b,c$ to agree. Therefore a progression of ideals in $\mathcal E_S$
would have zero valuation step outside $S$. Inside $S$ it would give
$u,u+h,u+2h\in D_S$ with $h\ge0$ and $h\ne0$.
If the final vector has a zero coordinate, all three vectors have a zero
there and lie in $A^k$, a contradiction. Otherwise the final vector is
in the new cube. At every moving coordinate its three entries must be
$0,1,2$. The initial vector consequently lies in $A^k$, so every fixed
coordinate belongs to $A\cap\{1,2\}=\{1\}$. The middle vector is the
deleted all-ones vector. This proves progression avoidance.

\paragraph{The density calculation.}
Set
\[
 w_{\mathfrak p}=(\operatorname N\mathfrak p)^{-1},\qquad
 T(w)=\sum_{\substack{a\in A\\a\ge1}}w^{a-1},\qquad
 p_{\mathfrak p}=(1-w_{\mathfrak p})\sum_{a\in A}w_{\mathfrak p}^{a},
 \qquad C_S=\prod_{\mathfrak p\notin S}p_{\mathfrak p}.
\]
All these products converge, and $C_S>0$. We use the classical ideal
counting asymptotic
$\#\{I:\operatorname NI\le X\}\sim\kappa_KX$, $\kappa_K>0$;
see \href{https://doi.org/10.1016/j.exmath.2006.08.001}{Murty--Van Order,
ideal-counting theorem}. Divisibility by an ideal $J$ has limiting proportion
$(\operatorname NJ)^{-1}$. Finite inclusion--exclusion gives independent
geometric valuation laws at any fixed finite set of prime ideals.
The unbounded conditions are justified by truncation: failure of the
greedy rule at a sufficiently large prime ideal requires
$\mathfrak p^2\mid I$. The counting asymptotic also gives a uniform
bound $\#\{I:\operatorname NI\le X\}\ll_K X$, so the upper density of
the omitted square-divisibility events is
$O_K(\sum_{\operatorname N\mathfrak p>Y}(\operatorname N\mathfrak p)^{-2})$,
which tends to zero. Large valuations at finitely many remaining primes
have a uniformly vanishing geometric tail. Thus both densities exist,
and subtraction of the deleted positive orthant and added cube gives
\begin{equation}\label{eq:ideal-natural-delta}
 d(\mathcal E_S)-d(\mathcal G_K)
 =C_S\prod_{\mathfrak p\in S}(1-w_{\mathfrak p})w_{\mathfrak p}
 \left[\prod_{\mathfrak p\in S}(1+w_{\mathfrak p})-1
                    -\prod_{\mathfrak p\in S}T(w_{\mathfrak p})\right].
\end{equation}

\paragraph{Why some finite exchange is profitable.}
The support of $T$ is contained in $\{0\}\cup\{2,3,\ldots\}$.
Every integer at least two is $2a+3b$ for some $a,b\ge0$, whence
\[
 T(w)\le\frac1{(1-w^2)(1-w^3)},\qquad
 \prod_{\mathfrak p\in S}T(w_{\mathfrak p})
       \le\zeta_K(2)\zeta_K(3)
       \le\bigl(\zeta(2)\zeta(3)\bigr)^d<2^d.
\]
The middle inequality follows prime by prime: at most $d$ prime ideals
lie over a rational prime, and each has norm at least that prime.
For completely rational bounds, integral tail estimates give
\[
 \zeta(2)<\sum_{n=1}^{10}n^{-2}+\frac1{10}
    =\frac{2095337}{1270080}<\frac{33}{20},\qquad
 \zeta(3)<\sum_{n=1}^{4}n^{-3}+\frac1{32}
    =\frac{2089}{1728}<\frac{121}{100};
\]
their last two bounds multiply to $3993/2000<2$.
On the other hand, $\sum_{\mathfrak p}(\operatorname N\mathfrak p)^{-1}$
diverges: otherwise the Euler product at $1$ would converge, contrary
to the ideal-counting asymptotic. Thus
$\prod_{\mathfrak p\in S}(1+w_{\mathfrak p})$ is unbounded along
increasing finite sets $S$.

Consequently the finite rational test
\begin{equation}\label{eq:ideal-natural-certificate}
                  \prod_{\mathfrak p\in S}(1+w_{\mathfrak p})>1+2^d
\end{equation}
eventually succeeds. Moreover
$C_S\ge\zeta_K(2)^{-1}>2^{-d}$, because $0,1\in A$.
The expression obtained from the right side of
\eqref{eq:ideal-natural-delta} by replacing $C_S$ by $2^{-d}$ and
$\prod T(w_{\mathfrak p})$ by $2^d$ is therefore a positive rational
lower bound for the gain whenever \eqref{eq:ideal-natural-certificate}
holds. Enumerating and factoring rational primes gives a terminating
search for $S$; factoring an input ideal and testing its valuation
digits decides membership in $\mathcal E_S$. This proves the theorem.

\paragraph{A short certificate over $\mathbb Q(\sqrt{-23})$.}
Write $\omega=(1+\sqrt{-23})/2$, whose minimal polynomial is
$X^2-X+6$. At each of $p=2,3$, its two distinct roots modulo $p$ are
$0,1$. The four prime ideals $(p,\omega)$ and $(p,\omega-1)$ therefore
have norms $2,2,3,3$. Take these four ideals for $S$.
The sharper bound
\[
                T(w)\le\overline T(w):=1+w^2+w^3+\frac{w^8}{1-w}
\]
follows from $A\cap[1,8]=\{1,3,4\}$. Exact arithmetic gives
\[
 \overline T(1/2)=\frac{177}{128},\qquad
 \overline T(1/3)=\frac{5023}{4374},\qquad
 \prod_{\mathfrak p\in S}(1+w_{\mathfrak p})-1
       -\prod_{\mathfrak p\in S}\overline T(w_{\mathfrak p})
       =\frac{16658080679}{34828517376}.
\]
Also $\prod_{\mathfrak p\in S}(1-w_{\mathfrak p})w_{\mathfrak p}=1/324$.
Since $\prod_{\mathfrak p\in S}(1-w_{\mathfrak p}^2)=4/9$ and
$\zeta_K(2)\le\zeta(2)^2<(5/3)^2$, we have
\[
 C_S\ge\prod_{\mathfrak p\notin S}(1-w_{\mathfrak p}^2)
       =\frac{\zeta_K(2)^{-1}}{4/9}>\frac{81}{100}.
\]
The resulting explicit natural-density gain is
\begin{equation}\label{eq:ideal-natural-23-gain}
 d(\mathcal E_S)-d(\mathcal G_K)
 >\frac{16658080679}{13931406950400}>0.0011957.
\end{equation}
The supplied rational checker verifies the prime reductions and every
number in this certificate. It does not substitute for the all-fields
proof above.

\paragraph{The precise boundary.}
The theorem treats integral-ideal ratios. The displayed exchange permits
mixed-step triples: for two selected coordinates, $(0,3),(1,2),(2,1)$
all satisfy the new rule, and further coordinates may be fixed at one.
Accordingly this construction does not answer Best et al.'s Question~5.4
about a \emph{natural}-density improvement for fractional-ideal ratios.
It also does not determine the optimal natural density. Those conclusions
must be separated from a strict natural-versus-upper-density theorem.
